\documentclass[11pt,a4paper]{article}
\usepackage[utf8]{inputenc}
\usepackage[margin=2.3cm]{geometry}
\usepackage{parskip}
\usepackage[hidelinks]{hyperref}
\usepackage{enumitem}
\usepackage{xcolor}
\usepackage{longtable}
\usepackage{booktabs}
\setlist[itemize]{leftmargin=1.4em, topsep=2pt}

% The letter is bound at the front of the revised article. Its pages are numbered
% R1, R2, ... so that every "p. n" below refers unambiguously to a page of the
% manuscript itself, whose own numbering begins again at 1.
\renewcommand{\thepage}{R\arabic{page}}

\definecolor{qbg}{RGB}{243,243,238}
\newcommand{\rc}[2]{%
  \par\medskip\noindent
  \colorbox{qbg}{\parbox{\dimexpr\linewidth-2\fboxsep}{%
  \textbf{#1}\par\smallskip\itshape #2}}\par\nopagebreak}
\newcommand{\resp}{\par\medskip\noindent\textbf{Response.}\ }
\newcommand{\where}[1]{\par\smallskip\noindent\textcolor{blue!55!black}{\textbf{Where to find it:} #1}\par}

\title{\vspace{-2.4cm}Response to the Editor and Reviewers\\[6pt]
\large IJADS-312370\\[2pt]
\normalsize ``From Signal Fusion to Asset Allocation: A Decision-Theoretic Model for
Portfolio Construction Under Regime-Based Sentiment and Volatility''\\[4pt]
\normalsize Anandkumar Pardeshi and Sujata Deshmukh}
\author{}
\date{}

\begin{document}
\maketitle
\vspace{-2.4cm}

\noindent\rule{\linewidth}{0.4pt}
\smallskip

\noindent\textbf{How to read this document.} Each recommendation is quoted in full, in the
shaded box, followed by our response and by the exact location of the change. Every
recommendation is answered individually: the three instructions in the Editor's letter,
three from Reviewer~1, and seven mandatory plus five optional from Reviewer~2.
A consolidated index of all eighteen, with its location, is given on the final page.

\noindent\textbf{Page numbers.} This response is bound at the front of the revised
article, as the Editor requested. Its pages carry the prefix \textbf{R}. Every
``p.~$n$'' in the text below refers to page~$n$ of the \emph{manuscript}, whose own
numbering restarts at~1 after this response ends.

\smallskip
\noindent\rule{\linewidth}{0.4pt}

\section*{Summary of the revision}

We thank both reviewers and the Editor for a careful and demanding review. Every one of
Reviewer~1's three mandatory recommendations and Reviewer~2's seven mandatory
recommendations has been addressed, all five of Reviewer~2's optional suggestions have
been implemented, and the three instructions in the Editor's covering letter are answered
in the next section.

The section structure and headings of the submitted version are retained unchanged, so
that the revision can be read against the original section by section. Sections~1 and~2
are substantially the text that was submitted, corrected for the faults Reviewer~1 and
Reviewer~2 identified and extended with the new references. Sections~3 to~5 retain the
submitted headings, and the material under them has been rewritten where the evidence
required it. The one structural addition is Section~4.3, which holds the robustness,
holdout, validation, ablation and inference analyses the reviewers asked for; there was no
subsection in the submitted version under which that material could sit.

\subsection*{One change larger than a normal revision}

We describe this at the outset, because the reviewers will notice it immediately.

Reviewer~2's third recommendation asked for robustness checks across rolling windows and
for sensitivity analysis. Implementing those checks exposed a defect in our own
implementation. The fused signal was being passed into the mean--variance objective in the
raw units of the underlying indicators; on-balance volume takes values of order $10^{8}$
to $10^{9}$, whereas annualised variances are of order $10^{-2}$. The linear term
therefore dominated the quadratic risk term by roughly ten orders of magnitude, the risk
model had no influence on the solution, and the optimiser returned a corner solution
holding a single asset on almost every rebalance date. This is precisely the behaviour
Reviewer~2 identified in recommendation~4 as inconsistent with the described
risk-budgeting approach. The reviewer was right, and the cause was a defect rather than a
design choice.

A second defect surfaced from the same investigation. The regime-dependent volatility
targets were specified as absolute levels of 10 to 20~per cent annualised. On a
diversified portfolio those levels are never reached, so the constraint bound in one
rebalance out of 224, and the regime layer --- the central mechanism of the paper --- had
no effect on the portfolio at all.

We corrected both defects. Expected returns are now scaled to a return metric
(Equation~(4), p.~15), and the regime risk budget is expressed as a fraction of the
portfolio's own volatility rather than as an absolute target (Equation~(7), p.~16).
Following Reviewer~2's optional suggestion~4, we also extended the evaluation from the
original five-year, ten-equity study to an eighteen-year, ten-instrument multi-asset study
covering four major stress episodes. The evaluation window runs to every session available
at the date the analysis was run, 26~August 2026, rather than stopping at the original
submission date. This is not a favourable window chosen after the fact: the 2008--2023
block reproduces exactly the figures the corrected pipeline produced before the extension,
and the 665 sessions from January~2024 onward postdate every specification decision in the
study. We therefore report them separately, in a new Section~4.3.1, as an out-of-sample
holdout. On that holdout the framework records a higher Sharpe ratio and a smaller maximum
drawdown than all three benchmarks, including the two that beat it on Sharpe ratio before
2024.

The consequence is that all reported numbers have changed, and the conclusions are
narrower than in the submitted version. The framework's drawdown and stress-period results
survive and are now supported across 18 of 18 specifications; the claim of significant
alpha does not survive and has been removed; and the signal fusion layer is now reported as
a tested null, since ablation shows it contributes nothing. We judged that reporting
corrected results with weaker claims was the only defensible course, and we would rather
present this to the reviewers now than have it discovered later.

\bigskip
\noindent\rule{\linewidth}{1pt}
\section*{The Editor's instructions}

\rc{E.1}{To help the reviewers check that you have made the required corrections, please
insert at the beginning of your revised article a detailed response to the reviewers'
recommendations. Make sure you address each recommendation thoroughly and methodically.
You should also show where the reviewers can find your change in the revised article by
referring to the page number and preferably highlighting the updated text.}

\resp This response is bound at the front of the revised article, ahead of the title page,
and its pages are numbered separately so that no page reference is ambiguous. Each of the
fifteen reviewer recommendations is quoted in full and answered individually, and each
answer ends with the section, table, figure, equation and page number at which the change
can be found. The final page, p.~R13, carries a consolidated index of all eighteen items.

On highlighting: we have not applied colour highlighting to the manuscript body, for a
reason we would rather state than leave the reviewers to infer. Because the two defects
described above changed every reported number, Sections~3 to~5 are new almost in their
entirety, and a highlighted copy would be almost uniformly highlighted, which would convey
less than this index does. Sections~1 and~2 are the opposite case: they are close to the
submitted text, and the specific changes there are enumerated under R1.1, R1.2 and R2.6
below. If the Editor would prefer a change-marked copy in addition, we will supply one on
request.
\where{This document, pp.~R1--R13; index on p.~R13.}

\rc{E.2}{When revising your article please ensure that your reference list is up to date
and that any recent articles, including those from IJADS, which are of relevance to your
article are included.}

\resp The reference list has been expanded from 33 to 54 entries. Ten are from 2025 or
2026. Three are from this journal and are cited where they bear on the argument rather
than in a block: Khodamoradi et~al. (2023) on robust mean-conditional-value-at-risk
portfolio optimisation, in Section~2.4 where the responses to input instability are
compared; Zhang et~al. (2025) on the divergence between in-sample and out-of-sample model
selection, in Section~2.3 where the evaluation design is motivated; and George and
Murugesan (2025) on sentiment analysis for equities, in Section~2.2.

The remaining new entries are the works the corrected methodology now depends on and which
the submitted version did not cite: Markowitz (1952), Ledoit and Wolf (2004), Grinold and
Kahn (2000), Hamilton (1989), Newey and West (1987), Politis and Romano (1994), DeMiguel
et~al. (2009), Harvey and Liu (2015), López de Prado (2016), Moreira and Muir (2017),
Kritzman et~al. (2012), Baker and Wurgler (2007) and Loughran and McDonald (2011);
together with recent allocation literature: Asimit et~al. (2025), Ha et~al. (2026), Hood
and Raughtigan (2025), Li et~al. (2025) and Luo and Mulvey (2026).
\where{References, pp.~41--45; Sections~2.2, 2.3 and~2.4, pp.~6--9.}

\rc{E.3}{We look forward to receiving your resubmission within the next 30 days.}

\resp The revision was completed within the period allowed.

\bigskip
\noindent\rule{\linewidth}{1pt}
\section*{Reviewer 1}

\subsection*{Changes which must be made before publication}

\rc{R1.1}{Please fix your literature review. A typical academic paper should have around
45 academic references with about 10 relevant references published in ACADEMIC JOURNALS in
2025 and 2026. Remove all self-citations and limit them to no more than 2--3 references.
References should be recent, directly relevant to the study, complete with full
bibliographic information, and drawn primarily from scholarly journals in analytics and
related fields rather than conference proceedings, professional magazines, trade
publications, or other non-scholarly sources. Please avoid excessive citations to a single
journal, including our own, or to a single author. As a general guideline, no more than
three citations should be attributed to the same journal or author. In addition, please
avoid citing references in clusters of more than three at a time.}

\resp We take each requirement in turn, and report one on which we have not complied.

\begin{itemize}
\item \textbf{Number of references.} Expanded from 33 to \textbf{54}, against the
approximately 45 suggested.
\item \textbf{Ten recent journal articles.} \textbf{Ten} of the 54 were published in 2025
or 2026, which meets the requested figure exactly: Asimit et~al. (2025), George and
Murugesan (2025), Ha et~al. (2026), Hood and Raughtigan (2025), Li et~al. (2025), Lu
et~al. (2025), Luo and Mulvey (2026), Sami et~al. (2025), Shang (2025) and Zhang et~al.
(2025). All ten are articles in peer-reviewed journals.
\item \textbf{Self-citations.} There are none, and there were none in the submitted
version.
\item \textbf{Full bibliographic information.} Every new entry was resolved against
CrossRef, so authors, title, journal, volume, issue, pages and DOI are taken from the
publisher record rather than from memory, and every DOI resolves. Two entries carry less
than the full set, correctly: Grinold and Kahn (2000) is a monograph and carries edition
and publisher instead, and Ha et~al. (2026) appears in a journal that assigns article
numbers rather than page ranges.
\item \textbf{Scholarly sources.} All 54 entries are journal articles apart from the one
monograph just named, which is cited for the information-ratio scaling used in
Equation~(4). There are no conference proceedings, magazines, trade publications, working
papers or preprints.
\item \textbf{Author concentration.} No author appears more than twice.
\item \textbf{Citation clusters.} Every citation command in the manuscript was checked
programmatically. No citation groups more than \emph{two} references, which is inside the
limit of three.
\item \textbf{Citations to this journal.} Three, which is at the stated limit and not
above it.
\end{itemize}

\textbf{The limit of three references per journal: we have not complied, and we set the
reasoning out for the Editor's decision rather than leave it unexplained.} Twenty-four of
the 54 entries are from \emph{IEEE Access}. All 24 were in the submitted version. We
retained the submitted reference list in full rather than replacing part of it, because
those 24 articles are the specific prior work the study's framing rests upon --- the
signal-fusion and sentiment-integration literature from which the signal-allocation gap is
identified in Sections~1.2 and~2.3 --- and because Tables~1, 2 and~3 attribute particular
methodological claims to particular ones of them. Removing two-thirds of that literature
to satisfy a counting rule would have required us either to delete those three comparative
tables or to re-attribute their rows to works that do not make the claims in question.

Every other journal in the list is within the limit. The maximum of three is reached by
this journal, by The Journal of Portfolio Management and by The Journal of Finance; no
other journal supplies more than two.

We recognise that this leaves one requirement unmet, and we do not wish to appear to have
overlooked it. If the Editor would prefer the limit applied strictly, we will reduce
\emph{IEEE Access} to three entries and rebuild Sections~2.1 to~2.4 and Tables~1 to~3
around the remaining literature. We would rather be told which of the two the journal
prefers than make that judgement unilaterally.
\where{References, pp.~41--45; new citations throughout Sections~1 and~2, pp.~2--9;
Tables~1--3, pp.~6, 7 and~9.}

\rc{R1.2}{The paper has some serious acronym-related problems. You cannot assume the
reviewers or readers are familiar with the acronyms you have used in your paper. You need
to expand an acronym once in the abstract upon the first use and once in the body of your
manuscript upon the first use. Remember, acronym overuse reduces readability.}

\resp Both halves of this recommendation have been acted on: every acronym is now expanded
at first use, and the total number of acronyms has been reduced.

The abstract contains no unexpanded acronym. In the body, first use and expansion were
checked programmatically against the compiled manuscript, so the ordering is verified
rather than asserted:

\begin{itemize}
\item open-high-low-close-volume (OHLCV), Section~3.1.1, p.~10
\item relative strength index (RSI), moving average convergence divergence (MACD), money
flow index (MFI) and on-balance volume (OBV), Section~3.1.2, p.~11
\item Gaussian mixture model (GMM), Akaike information criterion (AIC) and Bayesian
information criterion (BIC), Section~3.1.4, p.~13
\item information coefficient (IC), Section~3.2.2, p.~15
\item compound annual growth rate (CAGR) and maximum drawdown (Max DD), Section~3.3.2,
p.~18
\item percentage points (pp), caption to Table~6, p.~23
\item large language models (LLMs), Section~1.1, p.~3
\item radial basis function (RBF), long short-term memory (LSTM) and gated recurrent unit
(GRU), Table~1, p.~6
\item machine learning (ML), support vector machines (SVM) and convolutional neural
network (CNN), Table~2, p.~7
\end{itemize}

On overuse: we now write ``maximum drawdown'' in running text and reserve the abbreviation
for table headings, and the technical indicators are named in full in Table~4 (p.~12),
which is where a reader meets them first. The misspelling of ``LLMs'' as ``LLPs'' in
Section~2.2 of the submitted version is corrected.
\where{Abstract, p.~1; Sections~1.1, 3.1.1, 3.1.2, 3.1.4, 3.2.2 and~3.3.2; Tables~1 and~2,
pp.~6--7; Table~4, p.~12; Table~6 caption, p.~23.}

\rc{R1.3}{Please carefully polish and clean up all figures and tables. All figures and
tables in a paper must be consistent. Please make sure all figures are high-resolution
(300 dpi).}

\resp All eight figures have been regenerated at exactly 300~dpi from the corrected
pipeline, in a single consistent visual style: one typeface, one colour palette, matched
line weights, matched axis treatment and matched figure widths. The framework diagram
(Figure~1) was redrawn from scratch, both for resolution and because the submitted version
depicted a news and text ingestion path that the implementation does not contain.

Two further faults in the submitted version have been corrected, both of which bear on the
consistency the reviewer asks for. All seven figures previously shared the same
\verb|\label|, so every cross-reference in the submitted text resolved to the same figure;
each figure now carries a unique label and is referenced individually at the point it is
discussed. All eighteen tables have been rebuilt to a single format, with consistent rules,
column alignment, caption style and significant figures.
\where{Figures~1--8, pp.~11, 22, 23, 24, 27, 27, 29 and~30; Tables~1--18 throughout.}

\bigskip
\noindent\rule{\linewidth}{1pt}
\section*{Reviewer 2}

\subsection*{Changes which must be made before publication}

\rc{R2.1}{The manuscript has good potential, but several important improvements are needed
before publication. The methodology should be explained more clearly so that other
researchers can understand and reproduce the work.}

\resp Section~3 has been rewritten with reproduction as its explicit objective, and now
runs from p.~10 to p.~19. It retains the three-phase structure and every heading of the
submitted version, and fills them with the detail that was missing. It states the study
universe instrument by instrument, the data source and adjustment, the download window and
the resulting evaluation period; gives the exact formula and rolling window for all nine
input features in a single table; specifies the mixture model and its refit protocol;
states the fusion rule and the confidence term as equations; gives the optimisation problem
with every constraint and every parameter value; and sets out the walk-forward protocol and
the evaluation design. Nothing in the framework is now described only in prose.

We also state explicitly what the framework does \emph{not} contain, since the submitted
version's Section~3 described components that were not in the implementation. Section~3.1.1
now records that no news, social media or macroeconomic stream is ingested, and
Section~3.3.3 records which parameters are fixed ex ante and never re-optimised.
\where{Section~3, pp.~10--19.}

\rc{R2.2}{More implementation details are needed, particularly regarding the data sources
used, sampling frequency, sentiment analysis model, feature extraction process, Gaussian
Mixture Model settings, and the optimization solver applied.}

\resp Each of the six items is now specified explicitly.

\begin{itemize}
\item \textbf{Data sources.} Daily open, high, low, close and volume series from Yahoo
Finance through the \texttt{yfinance} interface, adjusted for splits and dividends, for the
ten instruments named individually in Section~3.1.1 (p.~10).
\item \textbf{Sampling frequency.} Daily. Download window 1~January 2007 to 27~August 2026;
evaluation period January~2008 to August~2026, comprising 4{,}689 sessions, of which the
final 665 are held out. Rebalancing is every 21 sessions (Sections~3.1.1 and~3.3.1).
\item \textbf{Sentiment analysis model.} Market-derived rather than textual, and now
specified under the submitted version's own heading (Section~3.1.3, p.~12): the equally
weighted mean of a five-session return, turnover relative to its own recent average,
realised volatility and volume instability, the last two entered with a negative sign. No
transformer and no lexicon are involved, because no text is read. See R2.o3 below, and
Section~4.3.3 (pp.~33--34), which validates that construction against four independently
published sentiment and stress series.
\item \textbf{Feature extraction process.} Table~4 (p.~12) gives the closed-form
definition and rolling window of all nine features, followed by the cross-sectional
standardisation in Equation~(1) and the definition of the two composites. Section~3.1.2
also names the indicators the framework does \emph{not} use.
\item \textbf{Gaussian mixture model settings.} Three components, full covariance matrices,
three random restarts, fixed seed, fitted on two standardised features, refitted at every
rebalance on an expanding window with a minimum of 252 observations (Section~3.1.4, p.~13).
\item \textbf{Optimisation solver.} CVXPY with the Clarabel interior-point solver;
covariance by Ledoit--Wolf shrinkage on 252 trailing sessions (Section~3.2.3, p.~16).
\end{itemize}

We also documented one implementation detail that is easy to miss and that we ourselves
initially handled incorrectly: mixture component indices permute freely between refits, so
labels must be anchored --- here by sorting components on fitted volatility --- or the
regime series is silently scrambled without raising an error (Section~3.1.4, p.~13).
\where{Section~3, pp.~10--19; Table~4, p.~12.}

\rc{R2.3}{The experimental results also need stronger statistical support. Although
improvements in Sharpe ratio and drawdown are reported, the manuscript would benefit from
including statistical significance tests, confidence intervals, robustness checks using
different rolling windows, turnover statistics, and sensitivity analysis to strengthen the
reliability of the findings.}

\resp All five are now reported, and this recommendation is what led to the corrections
described in the summary above.

\begin{itemize}
\item \textbf{Statistical significance tests.} Paired $t$-test and a Newey--West
$t$-statistic with ten lags on daily excess returns (Table~17, p.~35).
\item \textbf{Confidence intervals.} A stationary block bootstrap with 21-session blocks
and 2{,}000 replications gives interval estimates for the Sharpe ratio of the strategy, of
the benchmark, and of their difference (Table~17, p.~35).
\item \textbf{Robustness across rolling windows.} A full sweep of three random seeds
$\times$ two covariance estimation windows $\times$ three rebalancing intervals. All
eighteen cells are reported individually, not summarised (Table~12, p.~32).
\item \textbf{Turnover statistics.} Annual one-way turnover, turnover per rebalance, mean
and minimum effective number of assets, mean largest weight, and mean and minimum invested
exposure (Table~8, p.~26).
\item \textbf{Sensitivity analysis.} Transaction costs varied from 0 to 25 basis points
(Table~8, p.~26); the regime risk budget varied over five specifications (Table~16,
p.~35); component ablations at matched exposure (Table~15, p.~34).
\end{itemize}

We draw the reviewer's attention to what this analysis shows, because it does not favour
the submitted claim. The Sharpe improvement over the benchmark is \emph{not} statistically
significant: its bootstrap interval is $[-0.102, 0.574]$ and includes zero. Section~4.3.5
(p.~35) says so directly, and the abstract now states it too. What the evidence supports is
consistency rather than significance --- 18 of 18 specifications, 18 of 19 calendar years,
4 of 4 stress episodes, and a holdout period the specification never saw --- and the
manuscript is careful not to conflate the two.
\where{Sections~4.3.1--4.3.5, pp.~31--36; Tables~11, 12, 15, 16 and~17.}

\rc{R2.4}{The portfolio allocation results require further clarification. In Figure 5, the
portfolio appears concentrated in only a few assets during certain periods, which does not
fully align with the described risk-budgeting approach. This behaviour should be explained
more clearly.}

\resp The reviewer identified a real defect, and we are grateful for it. The concentration
was not a property of the strategy but an artefact of unit mismatch in the objective
function. In the submitted version the portfolio held exactly one asset on nearly every
date --- an effective number of assets of 1.00, with a single holding persisting for
1{,}364 of 1{,}421 sessions --- because the raw indicator scale overwhelmed the risk term
by roughly ten orders of magnitude. The reviewer's observation that this does not align
with the described risk-budgeting approach was exactly right: the risk model had no
influence on the solution at all.

Equation~(4) (p.~15) now scales the fused score to a return metric before it enters the
objective, which resolves the mismatch. The corrected allocation is diversified and stable:
the effective number of assets averages 5.40 and never falls below 4.60, the largest
position sits at the 25~per cent cap on average, and annual one-way turnover is 0.96 times
portfolio value. Section~4.2.4 (pp.~25--27) explains the failure mode explicitly rather
than quietly replacing the figure, since the diagnosis is itself a useful contribution and
a reader building a similar framework would benefit from it.
\where{Equation~(4), p.~15; Section~4.2.4, pp.~25--27; Table~8, p.~26; Figures~5 and~6, p.~26.}

\rc{R2.5}{The regime detection section also needs stronger justification. The choice of
three market regimes should be supported using model selection measures such as AIC, BIC,
or regime persistence statistics.}

\resp Table~7 (p.~24) now reports AIC and BIC for two through six components, together with
the persistence statistics for the three-state solution: self-transition probabilities,
mean run lengths, in-sample and out-of-sample state shares, and the annualised volatility
and momentum characterising each state.

We report an outcome that does not support our choice, and say so. Both AIC and BIC fall
monotonically across the range tested --- from 21944 and 22016 at two components to 20372
and 20599 at six --- so neither criterion selects three components. This is the expected
behaviour of a mixture fitted to persistent, heavy-tailed financial features, which absorbs
tail observations into additional components indefinitely and is rewarded for doing so.

The choice of three is therefore justified on the second of the two grounds the reviewer
offers, persistence, together with interpretability, and Table~7 supplies the evidence. The
three states are economically distinct: annualised volatility of 12.22, 20.04 and
47.07~per cent, with annualised momentum of $+16.04$, $-1.25$ and $-15.97$~per cent. They
persist: self-transition probabilities of 0.910 to 0.961, and mean run lengths of 11.1 to
25.3 sessions, which is long enough to act on at a monthly rebalancing frequency. The
transition matrix is close to tridiagonal --- the estimated probability of moving directly
between the calm and the stress state, in either direction, is zero --- so the volatility
ordering recovered by the model is economically meaningful rather than an artefact of
labelling. Finer partitions produce states too short-lived for the rebalancing schedule to
exploit.

Section~4.3.3 adds an independent check the reviewer did not ask for but which bears on the
same question: the identified states separate monotonically on a sentiment composite that
takes no part in fitting them, at $F = 19.34$, $p = 1.8 \times 10^{-8}$.
\where{Section~4.2.3, pp.~23--25; Table~7, p.~24; Figure~4, p.~24; Section~4.3.3,
pp.~33--34.}

\rc{R2.6}{In addition, the language and presentation require careful revision. There are
grammatical inconsistencies, sudden changes in writing style, and some unfinished
formatting elements, including incomplete AI declaration text.}

\resp All three faults are addressed.

\begin{itemize}
\item \textbf{Grammatical inconsistencies.} The manuscript has been given a full language
pass. Specific errors the reviewer will have noticed are corrected, including the
misspelling of ``LLMs'' as ``LLPs'' in Section~2.2, the malformed sentence opening
``A advanced regime detection framework'' in Section~1.5, and the non-word
``near-monosaccharic'' in Section~1.2.
\item \textbf{Changes in writing style.} The manuscript is now written in one voice
throughout, and a stray sentence of reviewer-style commentary that had been left inside
Section~3.1.1 of the submitted version has been removed.
\item \textbf{Unfinished formatting elements.} The AI declaration placeholder has been
replaced with a complete declaration naming the tools used and the use they were put to,
which was language editing only. We have also added acknowledgements, an author
contributions statement, a conflict of interest statement and a data and code availability
statement.
\end{itemize}
\where{Throughout; declarations on p.~40.}

\rc{R2.7}{Finally, the equations should be explained in greater detail, with clear
definitions of all parameters and constraints used in the optimization model.}

\resp Section~3.2.3 (pp.~16--17) now presents the optimisation as two explicit steps and
defines every symbol at the point of use. The objective terms and every constraint are
given, along with the value used and the reason for it: risk aversion $\lambda = 2$;
position cap $w_{\max} = 0.25$; non-negativity, which makes the strategy long-only;
turnover penalty $\kappa = c \cdot 252 / h$ with $c = 10$ basis points and $h = 21$
sessions, annualised so that trading cost is expressed on the same basis as return and
risk; the full-investment equality $\mathbf{1}^{\top} w = 1$, which is what makes the first
step purely relative; and the regime multiplier $k_r$, which sets total risk in the second
step.

Equations~(1) through~(7) each carry an explanation of what the expression does and why it
takes that form, including why cross-sectional rather than time-series standardisation is
used in Equation~(1), why Ledoit--Wolf shrinkage is required rather than the sample
covariance, why the information-ratio scaling of Equation~(4) is necessary rather than
cosmetic, and why the risk budget of Equation~(7) is relative rather than absolute.
\where{Section~3.2.3, pp.~16--17; Equations~(1)--(7), pp.~12--16.}

\subsection*{Suggestions which would improve the quality of the article}

All five have been implemented.

\rc{R2.o1}{The literature review can be strengthened by adding sharper critical comparison
between the proposed framework and recent portfolio allocation models.}

\resp Section~2.4 now compares the framework against nine allocation paradigms:
mean--variance optimisation, naive diversification, shrinkage-based optimisation,
signal-based allocation, risk-based construction, volatility management, regime-switching
allocation, reinforcement and learning-based allocation, and multimodal fusion. Table~3
(p.~9), which retains the submitted version's heading, sets out each paradigm's basis of
allocation and its principal limitation, and places this study among them --- with our own
limitation stated in the same terms as everyone else's, namely that the signal layer was
tested and found not to contribute.

Section~2.3 adds a discussion of evaluation practice, drawing on Harvey and Liu (2015),
López de Prado (2016) and Zhang et~al. (2025), which is what motivates the specification
sweep and the holdout rather than leaving those choices unexplained.
\where{Section~2, pp.~5--9; Table~3, p.~9.}

\rc{R2.o2}{Figures would benefit from more technical interpretation, especially regime
transitions, signal confidence behaviour, allocation stability.}

\resp Every figure is now interpreted in the text rather than merely described, under the
submitted version's own per-figure headings, and the three topics named are each treated
directly.

\begin{itemize}
\item \textbf{Regime transitions.} Analysed through the transition matrix and run lengths
in Section~4.2.3, including the finding that direct calm-to-stress transitions have
estimated probability zero, and the per-state realised performance that shows where the
framework's advantage arises.
\item \textbf{Signal confidence behaviour.} Figure~7 (p.~28) shows the amplitude of each
composite over time, and Section~4.2.5 interprets it together with the information
coefficients in Table~9 (p.~28) --- including the point that variation in amplitude is a
property of the signal and not evidence that the signal is informative.
\item \textbf{Allocation stability.} Quantified in Table~8 (p.~26) --- effective number of
assets, largest weight, turnover --- and a new figure, Figure~6 (p.~26), plots invested
exposure with the transitional and stress states shaded, so that the de-risking mechanism
is directly visible rather than inferred.
\end{itemize}
\where{Sections~4.2.1--4.2.6, pp.~21--31; Figures~2--8; Tables~7, 8 and~9.}

\rc{R2.o3}{The sentiment-analysis section may be improved by explicitly stating whether
domain-specific transformer models or lexicon-based methods were used.}

\resp Neither, and Section~3.1.3 (p.~12) --- which keeps the submitted version's heading,
``Sentiment Analysis Pipeline'' --- has been rewritten to say so without ambiguity. No news
article, social media post, corporate filing, transcript or analyst note is read at any
point in the study, so neither a domain-specific transformer nor a sentiment lexicon is
applied: there is no text for either to act upon, and there is no named entity recognition
stage because there are no entities to recognise. The reason is a data constraint rather
than a preference: textual measurement needs a time-stamped corpus covering the traded
instruments, and a fund holding several hundred securities has no news flow of its own.

What the submitted version called a sentiment pipeline is a composite of four price and
volume features, now defined term by term in Table~4 and placed in the literature to which
it actually belongs --- the market-derived tradition of Baker and Wurgler (2007), in which
turnover and dispersion stand in for the survey. That reference has been added, and
Section~2.2 now carries a paragraph and a table row on this family of measures.

Naming a proxy, however, obliges one to show that it behaves like the thing it stands in
for, so the revision does not stop at disclosure. Section~4.3.3 (pp.~33--34) is new. It
aggregates the composite to the market level and correlates it against four published
series that enter the framework nowhere: the University of Michigan survey of consumer
sentiment, the Cboe volatility index, the St.~Louis Fed financial stress index and the
Chicago Fed national financial conditions index. All four carry the expected sign, all four
are significant in first differences, and the sign is correct in all sixteen
series-by-block combinations when the sample is cut into four periods (Table~13, p.~33).
The survey is the demanding case, because it contains no market data whatever; monthly
changes in it move with monthly changes in the composite at $+0.240$ ($t = +3.54$). The
identified market states also separate on the composite, monotonically, at $F = 19.34$,
$p = 1.8 \times 10^{-8}$ (Table~14, p.~33), although the composite takes no part in fitting
them --- which is what entitles the manuscript to call those states sentiment and
volatility states rather than volatility states alone.

The unfavourable half is reported just as fully. Section~4.2.5 (pp.~27--29) measures the
predictive content of the composite and finds none: its information coefficient against
next-session cross-sectional returns is $-0.026$ ($t = -2.26$), and the regime-conditional
weighting hypothesis --- that sentiment information becomes relatively more informative
under stress --- is not supported, the calm-to-stress contrast being insignificant
($t = -1.46$, $p = 0.145$) with point estimates running the other way. Splitting at the
holdout boundary shows the negative coefficient to be era-specific and to disappear after
2023, so the layer cannot be rescued by inverting its sign either, and the portfolio-level
ablation shows that removing it does not harm performance. A measure can be valid and still
be unprofitable. The two findings are kept apart in the manuscript, and both are stated.
\where{Section~3.1.3, p.~12; Section~2.2, pp.~6--7; Sections~4.2.5 and~4.3.3, pp.~27--29
and~34--35; Tables~9, 13, 14 and~15.}

\rc{R2.o4}{Future work may include cross-market validation, larger asset universes,
reinforcement learning-based adaptive parameter tuning.}

\resp Two of the three have been brought into the paper rather than deferred.

\textbf{Larger asset universes and cross-market validation.} The evaluation has moved to a
ten-instrument multi-asset universe spanning US, developed and emerging equity, government
and corporate fixed income, gold and real estate, over eighteen years (Section~3.1.1,
p.~10). The original mega-capitalisation equity universe is retained as a secondary study in
Section~4.3.6 (pp.~36--37), so the two are validated side by side. That secondary result is
unfavourable and is reported as such: the framework loses on risk-adjusted return in a
concentrated single-sector universe, 1.145 against 1.371, while retaining its drawdown
advantage, $-20.45$ against $-33.63$ per cent. Rather than omit it, we use it to identify
the condition under which the framework is useful --- a universe with genuine dispersion in
asset volatility --- which is a more informative result than a second favourable one would
have been.

\textbf{Reinforcement-learning parameter tuning.} Retained as future work, but now stated
concretely as learning the risk-budget multipliers rather than as a general aspiration,
with the overfitting caveats that apply to doing so (Section~4.4.2, p.~39). We have also
added two further extensions in the same list: textual sentiment applied to a universe
where an instrument-level news corpus exists, and a regime model that represents state
persistence explicitly.
\where{Section~3.1.1, p.~10; Section~4.3.6, pp.~36--37; Table~18, p.~36; Section~4.4.2,
pp.~38--39.}

\rc{R2.o5}{The manuscript would benefit from minor restructuring to reduce repetition in
Sections 4.2 and 4.3.}

\resp The repetition has been removed, and the submitted version's headings are kept while
doing it. In the submitted version, Section~4.2 (``Detailed Analysis of Visualizations and
Performance Metrics'') and Section~4.3 (``Theoretical and Practical Implications'') each
restated the other's material, with the regime mechanism explained three times over.

In the revision, Section~4.2 holds one subsection per figure, exactly as in the submitted
version, with each answering a distinct question and none repeating another. The
robustness, holdout, external validation, ablation, inference and secondary-universe
analyses the reviewers asked for are gathered into a single new Section~4.3, rather than
interleaved with the figure discussion where they would have duplicated it. Interpretation
that previously appeared twice is consolidated into Section~4.4. The result is that
Section~4 has four subsections where the submitted version's structure would have required
fourteen.
\where{Section~4, pp.~20--39.}

\bigskip
\noindent\rule{\linewidth}{1pt}
\section*{Changes not requested by the reviewers}

For completeness, two changes were made that no reviewer asked for.

\begin{itemize}
\item \textbf{Abstract.} The claim of ``significant alpha'' has been removed, because the
corrected results do not support it --- the excess return is negative and insignificant
($p = 0.171$). The abstract now states the return sacrifice and the absence of statistical
significance alongside the positive findings. Its remaining wording is that of the
submitted abstract.
\item \textbf{Contributions.} The list in Section~1.5 has been extended with what the
evidence supports: the two-step separation of composition from total risk, the
matched-exposure ablation design, the out-of-sample holdout, and two documented failure
modes.
\end{itemize}

We believe the revised manuscript is a substantially more careful piece of work than the
version submitted, and we thank the reviewers for the comments that forced it.

\clearpage
\section*{Index of changes}

\noindent Page numbers refer to the manuscript, which begins after this response.

\begin{longtable}{@{}p{1.3cm}p{6.1cm}p{8.4cm}@{}}
\toprule
\textbf{Item} & \textbf{Recommendation} & \textbf{Where the change is} \\
\midrule
\endhead
E.1 & Response at front, page references & This document, pp.~R1--R9 \\
E.2 & Reference list up to date, IJADS articles & References pp.~41--45; Sections~2.2--2.4 \\
E.3 & Resubmit within 30 days & Completed within the period \\
\addlinespace
R1.1 & Literature review, 45 refs, 10 recent & References pp.~41--45; 54 refs, 10 from 2025--26. \emph{One point not complied with; see p.~R4} \\
R1.2 & Expand acronyms at first use & Abstract p.~1; Sections~1.1, 3.1.1, 3.1.2, 3.1.4, 3.2.2, 3.3.2; Tables~1--2 \\
R1.3 & Figures and tables, 300 dpi & Figures~1--8; Tables~1--18 \\
\addlinespace
R2.1 & Methodology reproducible & Section~3, pp.~10--19 \\
R2.2 & Implementation details & Section~3, pp.~10--19; Table~4, p.~12 \\
R2.3 & Statistical support & Sections~4.3.1--4.3.5; Tables~11, 12, 15, 16, 17 \\
R2.4 & Figure~5 concentration & Equation~(4) p.~15; Section~4.2.4, pp.~25--27 \\
R2.5 & Regime choice justified & Section~4.2.3, pp.~23--25; Table~7, p.~24 \\
R2.6 & Language, AI declaration & Throughout; declarations p.~40 \\
R2.7 & Equations explained & Section~3.2.3, pp.~16--17; Eqs.~(1)--(7) \\
\addlinespace
R2.o1 & Critical comparison in review & Section~2.4, p.~8; Table~3, p.~9 \\
R2.o2 & Technical interpretation of figures & Sections~4.2.1--4.2.6, pp.~21--31 \\
R2.o3 & Transformer or lexicon stated & Section~3.1.3, p.~12; Section~4.3.3, pp.~33--34 \\
R2.o4 & Cross-market, larger universe, RL & Section~3.1.1 p.~10; Section~4.3.6 pp.~36--37; Section~4.4.2 p.~39 \\
R2.o5 & Reduce repetition in \S4.2 and \S4.3 & Section~4, pp.~20--39 \\
\bottomrule
\end{longtable}

\end{document}
